\documentclass{article}
\usepackage[backend=biber,natbib=true,style=alphabetic,maxbibnames=50]{biblatex}
\addbibresource{/home/nqbh/reference/bib.bib}
\usepackage[utf8]{vietnam}
\usepackage{tocloft}
\renewcommand{\cftsecleader}{\cftdotfill{\cftdotsep}}
\usepackage[colorlinks=true,linkcolor=blue,urlcolor=red,citecolor=magenta]{hyperref}
\usepackage{amsmath,amssymb,amsthm,enumitem,float,graphicx,mathtools,tikz}
\usetikzlibrary{angles,calc,intersections,matrix,patterns,quotes,shadings}
\allowdisplaybreaks
\newtheorem{assumption}{Assumption}
\newtheorem{baitoan}{Bài toán}
\newtheorem{cauhoi}{Câu hỏi}
\newtheorem{conjecture}{Conjecture}
\newtheorem{corollary}{Corollary}
\newtheorem{dangtoan}{Dạng toán}
\newtheorem{definition}{Definition}
\newtheorem{dinhly}{Định lý}
\newtheorem{dinhnghia}{Định nghĩa}
\newtheorem{example}{Example}
\newtheorem{ghichu}{Ghi chú}
\newtheorem{hequa}{Hệ quả}
\newtheorem{hypothesis}{Hypothesis}
\newtheorem{lemma}{Lemma}
\newtheorem{luuy}{Lưu ý}
\newtheorem{nhanxet}{Nhận xét}
\newtheorem{notation}{Notation}
\newtheorem{note}{Note}
\newtheorem{principle}{Principle}
\newtheorem{problem}{Problem}
\newtheorem{proposition}{Proposition}
\newtheorem{question}{Question}
\newtheorem{remark}{Remark}
\newtheorem{theorem}{Theorem}
\newtheorem{vidu}{Ví dụ}
\usepackage[margin=2cm]{geometry}
\def\labelitemii{$\circ$}
\DeclareRobustCommand{\divby}{%
	\mathrel{\vbox{\baselineskip.65ex\lineskiplimit0pt\hbox{.}\hbox{.}\hbox{.}}}%
}
\setlist[itemize]{leftmargin=*}
\setlist[enumerate]{leftmargin=*}

\title{Đề Thi Cuối Kỳ Tổ Hợp {\it\&} Lý Thuyết Đồ Thị Hè 2026\\Final Exam: Combinatorics {\it\&} Graph Theory Summer 2026}
\author{Giảng viên ra đề: Nguyễn Quản Bá Hồng}
\date{\today}

\begin{document}
\maketitle
\begin{center}\large
	\emph{\textsf{Yêu cầu Bắt Buộc}}
\end{center}

\begin{enumerate}
	\item Được phép sử dụng tài liệu giấy (sách, báo, tập{\tt/}nháp viết tay, khăn giấy, giấy súc, etc.) không giới hạn.
	\item Cấm sử dụng thiết bị điện tử (trừ máy tính bỏ túi không có tích hợp AIs), AI tools (Gemini, ChatGPT, Grok, etc.). Nếu phát hiện 0 điểm ngay lần đầu tiên \& thu bài -- không có bất cứ cảnh cáo nào.
	\item Nếu làm không được, viết định nghĩa sẽ được +2 điểm. Trong mỗi bài, nếu không làm được ý trước thì vẫn có thể sử dụng kết quả các ý trước để làm ý sau của bài toán.
	\item Thời gian thi: 2.5 hours. Có thể nộp bài sớm nếu tự sinh viên muốn hoặc bị giảng viên bắt. Nộp bài sớm vẫn phải ngồi lại phòng thi, không được về sớm, để chiêm nghiệm về thái độ học tập, làm việc của bản thân \& đặc biệt về cuộc đời.
	\item Nếu phát hiện đề sai rồi sửa lại được đề \& làm đúng thì sẽ được cộng thêm nhiều điểm.
	\item Đề thi tổng cộng $\Sigma = 100$ điểm.
\end{enumerate}

\begin{center}\large
	\emph{\textsf{Đề Thi Chính Thức}}
\end{center}

%------------------------------------------------------------------------------%
\begin{baitoan}[Basic graphs -- Các đồ thị cơ bản]
	\begin{enumerate}[label=(\alph*)]
		\item \emph{(5 điểm)} Cho ví dụ \& vẽ đồ thị đơn, đa đồ thị, đồ thị tổng quát có $n$ đỉnh với mỗi $n\in[5]$. Viết cụ thể tập đỉnh $V$ \& tập $E$ cho từng ví dụ.
		\item \emph{(10 điểm)} Với đồ thị đơn $n\in\mathbb{N}^\star$ đỉnh tổng quát, tính số cạnh lớn nhất có thể. Viết $|V|,|E|$, $\{\deg v;v\in V\}$ trong trường hợp cực đại đó. Giải thích tại sao không thể xác định được với đa đồ thị hay đồ thị tổng quát?
	\end{enumerate}
\end{baitoan}

\begin{proof}[Lời giải toán học nghiêm ngặt]
	\emph{(a)} Xét ví dụ đại diện với $n = 3$, đặt tập đỉnh $V = \{1, 2, 3\}$. (Các trường hợp $n \in \{1, 2, 4, 5\}$ được xây dựng hoàn toàn tương tự).
	\begin{itemize}
		\item \emph{Đồ thị đơn (Simple graph):} Không chứa khuyên và cạnh bội. Tập cạnh ví dụ: $E = \{\{1,2\}, \{2,3\}, \{3,1\}\}$ (tương ứng với một hình tam giác).
		\item \emph{Đa đồ thị (Multigraph):} Không có khuyên nhưng cho phép cạnh bội. Tập cạnh ví dụ: $E = \{e_1, e_2, \{2,3\}, \{3,1\}\}$, trong đó $e_1, e_2$ cùng nối cặp đỉnh $\{1,2\}$.
		\item \emph{Đồ thị tổng quát (Pseudograph):} Cho phép cả cạnh bội và khuyên. Tập cạnh ví dụ: $E = \{e_1, e_2, \{2,3\}, \{3,1\}, e_3\}$, trong đó $e_3 = \{3,3\}$ là một khuyên tại đỉnh $3$.
	\end{itemize}
	\emph{(b)} Đối với một đồ thị đơn $G = (V,E)$ có $n$ đỉnh, số cạnh đạt cực đại khi và chỉ khi mọi cặp đỉnh phân biệt đều kề nhau. Khi đó $G$ đẳng cấu với đồ thị đầy đủ $K_n$.
	\begin{itemize}
		\item Số đỉnh: $|V| = n$.
		\item Số cạnh lớn nhất: $\max(|E|) = \binom{n}{2} = \frac{n(n-1)}{2}$.
		\item Tập các bậc của đỉnh: Do mọi đỉnh đều kề với toàn bộ $n-1$ đỉnh còn lại, ta có $\{\deg v \mid v \in V\} = \{n-1\}$.
	\end{itemize}
	\textit{Bản chất vô định của đa đồ thị và đồ thị tổng quát:} Trong hệ tiên đề của hai cấu trúc này, lực lượng của tập các cạnh song song giữa hai đỉnh bất kỳ (và tập các khuyên) là không bị giới hạn. Ta luôn có thể thêm một cạnh mới mà không vi phạm định nghĩa. Nói cách giải tích, hàm số đếm cạnh phân kỳ, tập $E$ không bị chặn trên ($\sup(|E|) = +\infty$). Do đó, việc tìm giá trị lớn nhất của $|E|$ là bài toán vô nghiệm.
\end{proof}

\begin{nhanxet}[Góc nhìn Khoa học Máy tính \& Chiêm nghiệm]
	Yêu cầu tìm ``số lượng cạnh lớn nhất'' của một đồ thị tổng quát là một phép thử tư duy cực kỳ sắc bén. Trong thế giới của Cấu trúc Dữ liệu, điều này ngô nghê y hệt như việc một Kỹ sư Kiểm thử phần mềm (\textsf{QA/Tester}) hỏi một \textsf{Intern}: \textit{``Số lượng Bugs tối đa em có thể tạo ra trong code là bao nhiêu?''} Câu trả lời, dĩ nhiên, là Vô cực. 
	
	Đồ thị tổng quát chính là hiện thân toán học của một vòng lặp vô hạn (\texttt{Infinite Loop}). Việc cố gắng lưu trữ một đa đồ thị có vô hạn cạnh bội vào Ma trận kề (\textit{Adjacency Matrix}) sẽ làm tràn biến \texttt{int}, tương tự như một cơn ác mộng làm tràn RAM và sinh ra lỗi \texttt{OutOfMemoryError}. Bài học rút ra: Đừng bao giờ cấp phát bộ nhớ tĩnh cho những thực thể vốn dĩ được sinh ra để chứa đựng sự vô hạn.
\end{nhanxet}

%------------------------------------------------------------------------------%
\begin{baitoan}[Some special graphs -- Vài đồ thị đặc biệt]
	\emph{(15 điểm)} Vẽ, viết ra $V,E,|V|,|E|$ cho đồ thị: (a) $P_{n-1}$. (b) Cycle $C_n$. (c) Complete graph $K_n$. (d) Complete bipartite graph $K_{m,n}$.
\end{baitoan}

\begin{proof}[Lời giải toán học nghiêm ngặt]
	Giả sử các đồ thị đều có nhãn chuẩn tắc.
	\begin{enumerate}[label=(\alph*)]
		\item \emph{Đường đi (Path) $P_{n-1}$:} Đồ thị có $n-1$ đỉnh.\\
		Tập đỉnh $V = \{v_1, v_2, \ldots, v_{n-1}\}$. Tập cạnh $E = \{v_i v_{i+1} \mid 1 \le i \le n-2\}$.\\
		Đại lượng: $|V| = n-1$, $|E| = n-2$.
		\item \emph{Chu trình (Cycle) $C_n$ ($n \ge 3$):}\\
		Tập đỉnh $V = \{v_1, v_2, \ldots, v_n\}$. Tập cạnh $E = \{v_i v_{i+1} \mid 1 \le i \le n-1\} \cup \{v_n v_1\}$.\\
		Đại lượng: $|V| = n$, $|E| = n$.
		\item \emph{Đồ thị đầy đủ (Complete graph) $K_n$:}\\
		Tập đỉnh $V = \{v_1, v_2, \ldots, v_n\}$. Tập cạnh $E = \{v_i v_j \mid 1 \le i < j \le n\}$.\\
		Đại lượng: $|V| = n$, $|E| = \binom{n}{2} = \frac{n(n-1)}{2}$.
		\item \emph{Đồ thị lưỡng phân đầy đủ (Complete bipartite graph) $K_{m,n}$:}\\
		Tập đỉnh phân hoạch thành 2 tập rời rạc $V = X \sqcup Y$ với $|X|=m, |Y|=n$ và $X \cap Y = \emptyset$.\\
		Tập cạnh $E = \{xy \mid x \in X, y \in Y\}$.\\
		Đại lượng: $|V| = m+n$, $|E| = mn$.
	\end{enumerate}
\end{proof}

\begin{nhanxet}[Mở rộng: Cấu trúc của Hệ thống Gợi ý - Recommendation Systems]
	Đồ thị lưỡng phân đầy đủ $K_{m,n}$ chính là kiến trúc dữ liệu lý tưởng của các hệ thống Recommendation như \textsf{Netflix} hay \textsf{Tinder}: Tập $X$ là \textit{Users} (Người dùng), Tập $Y$ là \textit{Items} (Phim ảnh/Đối tượng ghép đôi). Một đồ thị ``đầy đủ'' (\textit{complete}) đồng nghĩa với việc \emph{mọi} User đều tương tác với \emph{mọi} Item, tạo ra đúng $m \times n$ kết nối. Tuy nhiên, trên thực tế, các ma trận này luôn cực kỳ thưa thớt (\textit{Sparse Matrix}), và luôn tồn tại một lượng lớn các đỉnh có bậc bằng $0$ (\textit{isolated vertices}). Sự tồn tại của các đỉnh cô lập này chứng minh một định lý xót xa trong cuộc sống: Không phải lúc nào bạn cũng tìm được ``Match'' của đời mình!
\end{nhanxet}

%------------------------------------------------------------------------------%
\begin{baitoan}[Integer partition -- Phân hoạch số nguyên]
	\begin{enumerate}[label=(\alph*)]
		\item \emph{(5 đ)} Viết định nghĩa phân hoạch số nguyên của $n$, biểu đồ Ferrers, \& cho ví dụ. 
		\item \emph{(5 đ)} Viết định nghĩa phân hoạch liên hợp, biểu đồ Ferrer tương ứng, \& cho ví dụ.	
	\end{enumerate}
\end{baitoan}

\begin{proof}[Lời giải toán học nghiêm ngặt]
	\emph{(a) Phân hoạch số nguyên \& Biểu đồ Ferrers:}\\
	Một \emph{phân hoạch} của số nguyên dương $n$, ký hiệu $\lambda \vdash n$, là một cách biểu diễn $n$ dưới dạng tổng của một dãy các số nguyên dương giảm dần (hoặc không tăng) $\lambda = (\lambda_1, \lambda_2, \ldots, \lambda_k)$ thỏa mãn $\lambda_1 \ge \lambda_2 \ge \ldots \ge \lambda_k \ge 1$ và $\sum_{i=1}^k \lambda_i = n$. \\
	\emph{Biểu đồ Ferrers} là phương pháp biểu diễn hình học của $\lambda$. Ta vẽ các dấu chấm sắp xếp thành các hàng ngang và căn lề trái. Hàng thứ $i$ (từ trên xuống) chứa đúng $\lambda_i$ dấu chấm. \\
	\textit{Ví dụ:} Xét $\lambda = (4,2,1) \vdash 7$. Biểu đồ Ferrers có 3 hàng, số chấm lần lượt là 4, 2 và 1.
	
	\item(b) \textit{Phân hoạch liên hợp (Conjugate partition):}\\
	Phân hoạch liên hợp của $\lambda$, ký hiệu $\lambda'$, là phân hoạch thu được bằng cách đếm số lượng dấu chấm trên mỗi \emph{cột} của biểu đồ Ferrers của $\lambda$. Về mặt giải tích: $\lambda'_i = |\{j \mid \lambda_j \ge i\}|$. \\
	Biểu đồ Ferrers của $\lambda'$ thu được bằng phép chuyển vị (transpose) biểu đồ Ferrers gốc qua đường chéo chính (biến hàng thành cột, cột thành hàng). \\
	\textit{Ví dụ:} Chuyển vị biểu đồ $\lambda = (4,2,1)$, đọc số chấm theo cột ta thu được $\lambda' = (3,2,1,1) \vdash 7$.
\end{proof}

\begin{nhanxet}[Mở rộng: Quy hoạch động \& Đại số tuyến tính]
	Tại sao sinh viên Công nghệ phải học Phân hoạch số nguyên? Vì đây chính là bộ não toán học của bài toán \emph{Coin Change} (Đổi tiền) kinh điển trong \emph{Quy hoạch động (Dynamic Programming)}. Hơn nữa, phép toán lấy phân hoạch liên hợp $\lambda \mapsto \lambda'$ là một phép \emph{Involution} (tự nghịch đảo, tức $(\lambda')' = \lambda$), tương đương hoàn toàn với thuật toán chuyển vị một mảng 2D lởm chởm (\textit{Jagged Array}): \texttt{mat[i][j] = mat[j][i]} với độ phức tạp $\mathcal{O}(N \times M)$. Toán học chưa bao giờ tách rời khỏi thuật toán!
\end{nhanxet}

%------------------------------------------------------------------------------%
\begin{baitoan}[Pentagon number -- Số ngũ giác]
	\begin{enumerate}[label=(\alph*)]
		\item \emph{(5 đ)} Viết định nghĩa \& viết công thức toán của $|s(\lambda)|$ \& $|l(\lambda)|$.
		\item \emph{(5 đ)} Phát biểu định lý Euler cho số ngũ giác.
		\item \emph{(10 điểm)} Tính giá trị của hàm số $f(n)\coloneqq|{\cal E}(n)| - |{\cal O}(n)|$ với $n\in[10]$.
	\end{enumerate}
\end{baitoan}

\begin{proof}[Lời giải toán học nghiêm ngặt]
	Giả sử $\lambda = (\lambda_1, \lambda_2, \ldots, \lambda_m)$ là phân hoạch của $n$ thành các phần \emph{phân biệt} ($\lambda_1 > \lambda_2 > \ldots > \lambda_m \ge 1$).
	
	\emph{(a) Định nghĩa \& Công thức toán học:}
	Ký hiệu trong phép chứng minh song ánh của F. Franklin:
	\begin{itemize}
		\item \emph{Khối đáy $|s(\lambda)|$:} Là số lượng phần tử của phần nhỏ nhất (chiều dài của hàng cuối cùng trong biểu đồ Ferrers). Công thức: $|s(\lambda)| = \lambda_m$.
		\item \emph{Sườn dốc $|l(\lambda)|$:} Là số lượng lớn nhất các hàng liên tiếp (tính từ đỉnh) mà mỗi hàng giảm đúng 1 đơn vị. Công thức: $|l(\lambda)| = \max\{k \ge 1 \mid \lambda_i - \lambda_{i+1} = 1, \forall 1 \le i < k\}$.
	\end{itemize}
	
	\item(b) \textit{Định lý Euler về Số ngũ giác (Pentagonal Number Theorem):}\\
	Gọi ${\cal E}(n)$ và ${\cal O}(n)$ lần lượt là tập các phân hoạch của $n$ thành số lượng \emph{chẵn} và \emph{lẻ} các phần phân biệt. Chênh lệch phần tử giữa hai tập được cho bởi:
	\[
	|{\cal E}(n)| - |{\cal O}(n)| = \begin{cases} 
		(-1)^k & \text{nếu } n = \frac{k(3k \pm 1)}{2} \text{ với } k \in \mathbb{N}^\star \\ 
		0 & \text{trong các trường hợp còn lại (không phải số ngũ giác tổng quát).}
	\end{cases}
	\]
	
	\item(c) \textit{Tính $f(n) \coloneqq |{\cal E}(n)| - |{\cal O}(n)|$ cho $n \in [10]$:}\\
	Thay vì liệt kê thủ công toàn không gian trạng thái, ta thế trực tiếp $k$ vào công thức sinh số ngũ giác:
	\begin{itemize}
		\item $k=1 \implies n \in \{1, 2\} \implies f(1) = f(2) = (-1)^1 = -1$.
		\item $k=2 \implies n \in \{5, 7\} \implies f(5) = f(7) = (-1)^2 = 1$.
		\item $k=3 \implies n \in \{12, 15\}$, đã vượt ra khỏi khoảng $[10]$ (vòng lặp dừng).
	\end{itemize}
	Với mọi $n \in \{3,4,6,8,9,10\}$, do không mang dạng số ngũ giác tổng quát, ta kết luận ngay $f(n) = 0$.
\end{proof}

\begin{nhanxet}[Đòn bẩy $\mathcal{O}(1)$ của Toán học -- Time Complexity Trap]
	Bài toán này là một cái bẫy hoàn hảo về \textsf{Time Complexity}. Nếu sinh viên hồn nhiên lập trình hàm $f(n)$ bằng thuật toán quay lui (Backtracking) để liệt kê không gian trạng thái của $\mathcal{E}(n)$ và $\mathcal{O}(n)$, hệ thống sẽ phải gánh chịu thời gian chạy tiệm cận hàm mũ $\mathcal{O}(e^{C\sqrt{n}})$. Bằng việc hiểu sâu sắc Định lý Euler, ta có thể truy xuất giá trị bằng công thức toán như một \emph{Hash Function} với độ phức tạp $\mathcal{O}(1)$ và chốt hạ kết quả trong 10 giây. Làm kỹ sư giỏi không phải là người viết code "trâu" nhất, mà là người biết dùng Toán học để phá vỡ các giới hạn vật lý của CPU!
\end{nhanxet}

%------------------------------------------------------------------------------%
\begin{baitoan}[Regular graph -- Đồ thị chính quy]
	\emph{(10 đ)} Vẽ đồ thị $d$-chính quy có $7$ đỉnh với mỗi $d\in[6]$. Nếu không thể vẽ được về mặt toán học, chứng minh tại sao không vẽ được.
\end{baitoan}

\begin{proof}[Lời giải toán học nghiêm ngặt]
	Theo \emph{Định lý Bắt tay (Handshaking Lemma)}, trong một đồ thị vô hướng bất kỳ, tổng bậc của các đỉnh phải luôn bằng 2 lần số lượng cạnh: $\sum_{v\in V} \deg(v) = 2|E|$. Từ đó suy ra, tổng bậc luôn luôn phải là một số chẵn.
	
	Xét đồ thị $d$-chính quy có $n=7$ đỉnh. Bậc của mọi đỉnh đều là $d$, nên tổng bậc của đồ thị là $7d$. Để đồ thị tồn tại, $7d$ phải là số chẵn, kéo theo $d$ bắt buộc phải là số chẵn.
	\begin{itemize}
		\item \emph{Trường hợp không thể vẽ ($d \in \{1, 3, 5\}$):} Khi đó $7d \in \{7, 21, 35\}$ đều là các số lẻ. Điều này vi phạm nghiêm trọng Bổ đề Bắt tay. Ta kết luận hệ đồ thị này vô nghiệm về mặt toán học.
		\item \emph{Trường hợp có thể vẽ ($d \in \{2, 4, 6\}$):} 
		\begin{itemize}
			\item \emph{$d=2$:} Vẽ đồ thị chu trình $C_7$ (một đa giác đều 7 cạnh).
			\item \emph{$d=4$:} Vẽ $C_7$, sau đó từ mỗi đỉnh nối thêm 2 đường chéo sang 2 đỉnh cách nó đúng 1 bước (tạo thành một hình ngôi sao 7 cánh nội tiếp bên trong đa giác).
			\item \emph{$d=6$:} Bậc cực đại của đồ thị đơn 7 đỉnh là $7-1=6$. Vẽ đồ thị đầy đủ $K_7$ (mọi đỉnh nối chằng chịt trực tiếp với 6 đỉnh còn lại).
		\end{itemize}
	\end{itemize}
\end{proof}

\begin{nhanxet}[Quy tắc Sanity Check]
	Đây là một \emph{Unit Test} mẫu mực cho tư duy lập trình. Việc chạy một dòng rào cản sơ bộ: \texttt{if ((N * d) \% 2 != 0) return ERROR;} ngay từ đầu giúp giải phóng 50\% khối lượng công việc, thay vì lấy nháp ra hì hục phác thảo một cấu trúc đồ thị 3-chính quy 7 đỉnh không hề tồn tại trong vũ trụ. Data Validation luôn phải được thực thi trước khi phân bổ tài nguyên xử lý!
\end{nhanxet}

%------------------------------------------------------------------------------%
\begin{baitoan}[Euler's theorem{\tt/}handshaking lemma {\it\&} Havel--Hakimi algorithm -- Định lý Euler/bổ đề bắt tay {\it\&} thuật toán Havel--Hakimi]
	{\rm($\Sigma = 15$ điểm)}
	\begin{enumerate}[label=(\alph*)]
		\item {\rm(5 điểm)} Phát biểu định lý Euler \& định lý về thuật toán Havel--Hakimi. 2 định lý này áp dụng cho các loại đồ thị nào?
		\item {\rm(5 điểm) \cite[P10.1.13, p. 368]{Shahriari2022}} Chứng minh 1 dãy $d_1,d_2,\ldots,d_n$ là 1 graphic sequence khi \& chỉ khi dãy $n - d_n - 1,\ldots,n - d_2 - 1,n - d_1 - 1$ là graphic; áp dụng kiểm tra
		\begin{equation*}
			14,13,12,11,10,9,9,9,9,9,9,9,9,8,8,8,7,6,5,4,3,2,1
		\end{equation*}
		có phải là graphic sequence không.
		\item {\rm(5 điểm)} Sử dụng thuật toán Havel--Hakimi để kiểm tra dãy
		\begin{equation*}
			14,13,12,11,10,9,9,9,9,9,9,9,9,8,8,8,7,6,5,4,3,2,1
		\end{equation*}
		có phải là graphic sequence hay không.
	\end{enumerate}
\end{baitoan}

\begin{proof}[Lời giải toán học nghiêm ngặt]
	\emph{(a) Phát biểu định lý:}
	\begin{itemize}
		\item \textit{Định lý Euler (Bổ đề bắt tay):} Tổng bậc của các đỉnh trong đồ thị bằng hai lần số lượng cạnh ($\sum \deg(v) = 2|E|$). \textit{Phạm vi áp dụng:} Mọi loại đồ thị vô hướng.
		\item \textit{Định lý Havel-Hakimi:} Dãy số không âm giảm dần $D = (d_1, \ldots, d_n)$ là một graphic sequence (dãy bậc đồ thị) khi và chỉ khi dãy $D'$ (thu được bằng cách xóa $d_1$ và trừ $1$ đơn vị vào $d_1$ phần tử tiếp theo, rồi sắp xếp lại giảm dần) cũng là một graphic sequence. \textit{Phạm vi áp dụng:} Chỉ đồ thị đơn.
	\end{itemize}
	
	\item(b) \textit{Chứng minh bằng Đồ thị bù (Complement Graph):} 
	($\Rightarrow$) Giả sử dãy $D = (d_1 \ge \ldots \ge d_n)$ là graphic sequence, sinh ra đồ thị đơn $G$ có $n$ đỉnh. Trong đồ thị đơn, bậc lớn nhất của một đỉnh là $n-1$. Ta xét đồ thị bù $\overline{G}$ trên cùng tập đỉnh $V$ (hai đỉnh kề nhau trong $\overline{G} \iff$ chúng không kề nhau trong $G$). Do đó, bậc của đỉnh $v_i$ trong $\overline{G}$ chính xác bằng $(n-1) - \deg_G(v_i) = n-1-d_i$. Vì đồ thị phần bù của một đồ thị đơn vẫn là đồ thị đơn, tập các bậc này là một graphic sequence. Xếp dãy theo thứ tự không tăng, ta được $n-1-d_n, \ldots, n-1-d_1$. 
	Chiều ($\Leftarrow$) được chứng minh tương tự do phép toán lấy phần bù là tự nghịch đảo ($\overline{\overline{G}} = G$).
	
	\emph{Áp dụng kiểm tra:} Gọi dãy số đề bài gồm $23$ phần tử là $S$. 
	Tổng bậc của $S$ là: $\sum S = 14+13+12+11+10+(9 \times 8)+(8 \times 3)+7+6+5+4+3+2+1 = \mathbf{184}$. 
	Vì tổng bậc là một \emph{số chẵn}, nó hợp lệ qua màng lọc Parity Check! Xét dãy bù $S'$ ($d'_i = 22 - d_i$), tổng của nó là $\sum S' = 23 \times 22 - 184 = 322$ (vẫn chẵn). Dãy không thể bị loại, ta buộc phải dùng đệ quy Havel-Hakimi ở câu (c) để kết luận.
	
	\item(c) \textit{Khử đệ quy Havel-Hakimi:} 
	Do dãy số có nhiều phần tử lặp lại, ta ký hiệu $x^k$ là $k$ số $x$. Ta chạy thuật toán:
	\begin{align*}
		S_0 &= (14, 13, 12, 11, 10, 9^8, 8^3, 7, 6, 5, 4, 3, 2, 1) \\
		\xrightarrow{\text{Xóa } 14} S_1 &= (12, 11, 10, 9, 8^9, 7^3, 6, 5, 4, 3, 2, 1) \\
		\xrightarrow{\text{Xóa } 12} S_2 &= (10, 9, 8, 7^{12}, 6, 5, 4, 3, 2, 1) \\
		\xrightarrow{\text{Xóa } 10} S_3 &= (8, 7^5, 6^9, 5, 4, 3, 2, 1) \\
		\xrightarrow{\text{Xóa } 8} S_4 &= (6^{11}, 5^4, 4, 3, 2, 1) \\
		\xrightarrow{\text{Xóa } 6} S_5 &= (6^4, 5^{10}, 4, 3, 2, 1) \xrightarrow{\text{Xóa } 6} S_6 = (5^{10}, 4^4, 3, 2, 1) \\
		\xrightarrow{\text{Xóa } 5} S_7 &= (5^4, 4^9, 3, 2, 1) \xrightarrow{\text{Xóa } 5} S_8 = (4^{10}, 3^3, 2, 1) \\
		\xrightarrow{\text{Xóa } 4} S_9 &= (4^5, 3^7, 2, 1) \xrightarrow{\text{Xóa } 4} S_{10} = (3^{11}, 2, 1) \\
		\xrightarrow{\text{Xóa } 3} S_{11} &= (3^7, 2^4, 1) \xrightarrow{\text{Xóa } 3} S_{12} = (3^3, 2^7, 1) \xrightarrow{\text{Xóa } 3} S_{13} = (2^8, 1^2) \\
		\xrightarrow{\text{Xóa } 2} S_{14} &= (2^5, 1^4) \xrightarrow{\text{Xóa } 2} S_{15} = (2^2, 1^6) \xrightarrow{\text{Xóa } 2} S_{16} = (1^6, 0) \\
		\xrightarrow{\text{Xóa } 1} S_{17} &= (1^4, 0^2) \xrightarrow{\text{Xóa } 1} S_{18} = (1^2, 0^3) \xrightarrow{\text{Xóa } 1} S_{19} = (0^4)
	\end{align*}
	Thuật toán bào mòn dãy số một cách hoàn hảo và hội tụ về mảng toàn số $0$ hợp lệ. Không có phần tử âm nào sinh ra. Kết luận: Dãy $S$ \emph{LÀ MỘT GRAPHIC SEQUENCE}.
\end{proof}

\begin{nhanxet}[Đòn Tâm Lý Chiến Tuyệt Đỉnh -- Psychological Warfare]
	Đây chính là \emph{``Cạm bẫy lồng trong Cạm bẫy''} mang tầm vóc lịch sử! Ở các phiên bản trước của dãy số này, tận cùng là số \texttt{2} (độ dài 22), tổng bậc khi đó bằng $\mathbf{183}$ (Lẻ). Ai nhận ra sẽ kết luận ngay vô nghiệm bằng phép thử Parity Check $\mathcal{O}(1)$ cực kỳ ngầu. Nhưng ở đề này, giảng viên đã lẳng lặng \emph{thêm đúng số ``1''} vào cuối chuỗi. Sự thay đổi siêu nhỏ của "1 chiếc lá" đã bẻ tổng thành $\mathbf{184}$ (Chẵn), hồi sinh toàn bộ hệ thống từ trạng thái "Crash" sang "Graphic", vô hiệu hóa hoàn toàn mẹo tính nhanh. Sinh viên bị ném thẳng vào địa ngục đệ quy $\mathcal{O}(N^2)$ của thuật toán Havel-Hakimi thực thụ, phải ngồi trừ tay 19 vòng lặp. Sự tàn nhẫn và lãng mạn của Toán học hội tụ cả ở số \texttt{1} bé nhỏ này! Một bài học xương máu cho Dev: Sửa 1 dòng code, thay đổi toàn bộ kiến trúc!
\end{nhanxet}

%------------------------------------------------------------------------------%
\begin{baitoan}[Paths, trees {\it\&} degree sequences -- Đường đi, cây, {\it\&} dãy bậc]
	{\rm($\Sigma = 15$ điểm)} Trong khoa học máy tính, 1 đồ thị cây (tree) $T = (V, E)$ là 1 đồ thị đơn vô hướng, liên thông (connected) \& không chứa chu trình (cycle). Đã biết, mọi đồ thị cây có $n$ đỉnh thì luôn có đúng $|E| = n - 1$ cạnh.	
	\begin{enumerate}[label=(\alph*)]
		\item {\rm(5 điểm)} Viết định nghĩa của đường đi (path) \& chu trình (cycle) trong đồ thị.
		\item {\rm(5 điểm)} Kế thừa Định lý bắt tay (Handshaking lemma) ở Bài 6, hãy chứng minh: Mọi đồ thị cây có $n \ge 2$ đỉnh luôn tồn tại ít nhất $2$ đỉnh treo (đỉnh lá -- leaf, tức là đỉnh có bậc đúng bằng $1$).
		\item {\rm(5 điểm)} \emph{(Tư duy thuật toán)} Chứng minh: 1 dãy gồm $n$ số nguyên dương $d_1\ge d_2\ge\ldots\ge d_n\ge1$ ($n \ge 2$) là dãy bậc (degree sequence) của 1 đồ thị cây khi \& chỉ khi:
		\begin{equation*}
			\sum_{i = 1}^n d_i = 2n - 2.
		\end{equation*}
		{\sf Hint.} Chiều đảo $\Leftarrow$: Sử dụng quy nạp theo số đỉnh $n$. Nếu tổng các bậc bằng $2n - 2$ \& tất cả $d_i \ge 1$, chắc chắn phải có ít nhất 1 phần tử bằng 1. Thử ``ngắt'' đỉnh lá này đi \& nối vào 1 đồ thị cây $n-1$ đỉnh).
	\end{enumerate}
\end{baitoan}

\begin{proof}[Lời giải toán học nghiêm ngặt]
	\emph{(a) Định nghĩa:} 
	\textit{Đường đi (Path):} Là một dãy các đỉnh liên tiếp phân biệt $v_1, \ldots, v_k$ sao cho luôn tồn tại cạnh nối giữa $v_i$ và $v_{i+1}$. 
	\textit{Chu trình (Cycle):} Là một đường đi khép kín (đỉnh đầu trùng đỉnh cuối) có chiều dài $k \ge 3$.
	
	\item(b) \textit{Bắt tay để tìm Lá:} 
	Gọi $T$ là Cây $n \ge 2$ đỉnh, suy ra $|E| = n-1$. Tổng bậc đồ thị là $\sum \deg(v) = 2|E| = 2n-2$. Do tính liên thông, không đỉnh nào có bậc 0 (tức $\deg(v) \ge 1, \forall v$).
	Sử dụng phản chứng: Giả sử Cây có ít hơn 2 lá (tức có tối đa 1 đỉnh bậc 1, và $n-1$ đỉnh còn lại đều có bậc $\ge 2$). 
	Khi đó tổng bậc $\sum \deg(v) \ge 1(1) + (n-1)2 = 2n-1$. Điều này tạo mâu thuẫn trực tiếp vì $2n-1 > 2n-2$. Vậy Cây bắt buộc phải có ít nhất 2 đỉnh lá.
	
	\item(c) \textit{Thuật toán Cắt lá -- Leaf Plucking:}\\
	\textit{Chiều $\Rightarrow$:} Nếu đồ thị là Cây thì $|E| = n-1$, theo Bổ đề bắt tay tổng các bậc là $\sum d_i = 2(n-1) = 2n-2$. (Hiển nhiên)\\
	\textit{Chiều $\Leftarrow$ (Quy nạp toán học):} 
	\begin{itemize}
		\item \emph{Cơ sở:} $n=2$, dãy có tổng bằng $2(2)-2=2 \implies (1,1)$, đồ thị $K_2$ chính là Cây $P_2$. 
		\item \emph{Quy nạp:} Giả sử mệnh đề đúng với $n-1$. Xét dãy $n$ phần tử dương có tổng $2n-2$. Vì trung bình cộng $\frac{2n-2}{n} < 2$, phần tử nhỏ nhất bắt buộc $d_n = 1$. Lại có $2n-2 > n$ ($n \ge 3$), nên phần tử lớn nhất $d_1 \ge 2$.
		\item Ta "ngắt" chiếc lá này đi bằng cách xóa $d_n$ và trừ $1$ vào $d_1$. Dãy mới $D' = (d_1-1, d_2, \ldots, d_{n-1})$ có $n-1$ phần tử dương, tổng bằng $(2n-2)-1-1 = 2n-4 = 2(n-1)-2$. 
		\item Theo giả thiết quy nạp, $D'$ là dãy bậc của một Cây $T'$ có $n-1$ đỉnh. Phục hồi đồ thị gốc bằng cách thêm một đỉnh mới $v_n$ và nối duy nhất vào đỉnh tương ứng với $d_1-1$ trong $T'$. Đồ thị $T$ thu được liên thông, không sinh chu trình, và sở hữu chính xác dãy bậc ban đầu.
	\end{itemize}
\end{proof}

\begin{nhanxet}[Garbage Collection \& Pop Call Stack]
	Thuật toán \textit{``Ngắt lá''} (Leaf Plucking) ở Bài 7c chính là hiện thân của cơ chế \emph{Pop Call Stack} trong đệ quy (Recursion) và là cốt lõi của thuật toán mã hóa siêu việt \emph{Prüfer Code} (thuật toán nén cây $n$ đỉnh thành mảng $n-2$ phần tử). Trực giác này tương đương với thuật toán \emph{Topological Sort} và \emph{Garbage Collection} (Reference Counting) trong Hệ điều hành: Để tháo dỡ hoặc xây dựng một cấu trúc dữ liệu khổng lồ (DOM Tree) mà không bị Memory Leak, ta luôn bóc tách an toàn từ các node không có dependencies (lá bậc 1) và dồn ngược về Root!
\end{nhanxet}

%------------------------------------------------------------------------------%
\section*{Phụ lục: Hồ sơ Sinh hóa Thần kinh (Neuroendocrine Rollercoaster)}
\textit{Một đề thi tuyệt phẩm không sinh ra để kiểm tra trí nhớ, mà là một trải nghiệm UX (User Experience) thao túng trực tiếp não bộ sinh viên. Dưới góc nhìn Y sinh, 2.5 giờ trong phòng thi là một chuyến tàu lượn siêu tốc của hệ nội tiết:}
\begin{enumerate}
	\item \emph{Cú ``Hack'' Cortisol (Phút 1 - 5):} Khi bóc đề, \textit{Cortisol} (hormone căng thẳng) tăng vọt. Đọc đến quy chế \textit{``Được dùng giấy súc'', ``Nộp sớm ngồi lại chiêm nghiệm cuộc đời''}, sự mỉa mai đã bẻ gãy phản xạ ``Fight or Flight'', Cortisol rơi tự do và \textit{Endorphin} trào ra qua những tiếng cười khúc khích.
	\item \emph{Liệu pháp Serotonin (Phút 10 - 60):} Chép định nghĩa bài 1,2,3 nhờ quy chế Open-book. \textit{Serotonin} (hormone tự tôn) dâng cao. Cảm giác an toàn giả tạo xuất hiện: \textit{``Đề này mình qua môn chắc rồi!''}
	\item \emph{Cơn bão Dopamine \& Cú lừa Parity (Phút 60 - 90):} Chạm trán bài 6. Ai từng luyện đề năm cũ sẽ nghĩ tổng là lẻ. Khí thế bừng bừng cộng lại ra... 184 (Chẵn). Mắt trợn ngược. Cơn bão \textit{Adrenaline} bùng nổ khi nhận ra: Bắt buộc phải chạy đệ quy 19 bước bằng tay. Giảng viên đã out-play tất cả!
	\item \emph{Dòng chảy Acetylcholine (Phút 90 - 120):} Vượt qua Bài 7 bằng đệ quy cắt lá. Sóng não từ Beta chuyển sang Alpha. \textit{Acetylcholine} tiết ra liên tục, tạo \textit{Flow State} tuyệt đối để hoàn thiện chứng minh.
	\item \emph{Trạng thái Oxytocin (30 Phút cuối):} Nộp bài nhưng bị bắt ngồi lại. \textit{GABA} làm dịu hệ thần kinh. Sinh viên nhìn bạn bè vã mồ hôi đếm thủ công bài 6 và thực sự \emph{\textit{bắt đầu chiêm nghiệm cuộc đời}}. Sự kính trọng (Respect) dành cho người ra đề được mã hóa vĩnh viễn vào liên kết Nơ-ron!
\end{enumerate}

\printbibliography[heading=bibintoc]

\end{document}